\chapter{System Implementation Details}
\label{ch:implementation}

\section{Backend Services Architecture (Fastify 5 + TypeScript)}
The core backend application is built using **Fastify 5** and TypeScript, encapsulated inside the \texttt{apps/api} directory. Fastify was selected for its minimal runtime overhead, integrated schema-based request validation (via TypeBox and Ajv), and native asynchronous lifecycle hooks.

\subsection{Modular Directory Structure}
The backend enforces a clean modular architecture where each functional domain encapsulates its routes, services, schemas, and tests:
\begin{itemize}
    \item \texttt{src/modules/auth/}: User registration, login, JWT issuance, Google OAuth2 SSO exchange, and password reset flows.
    \item \texttt{src/modules/applications/}: Kanban CRUD operations, state transition dispatchers, stage lock controllers, and status event histories.
    \item \texttt{src/modules/jobs/}: Company career portal ingestion, multi-ATS synchronization, and multi-tier deduplication execution.
    \item \texttt{src/modules/mailbox/}: Gmail webhook receivers, OAuth credential encryption, message queueing, and recruiter review workflows.
    \item \texttt{src/modules/ai/}: HTTP RPC proxy communicating with the Python AI microservice for embeddings, resume fit scoring, and bullet point generation.
    \item \texttt{src/modules/queue/}: Transactional outbox queue runner and background task dispatchers.
\end{itemize}

\subsection{Empty Body Parsing and Fastify Middleware}
To eliminate standard Fastify HTTP 400 parsing errors when client fetch requests issue POST/DELETE calls with \texttt{Content-Type: application/json} but empty bodies, a custom content-type parser is registered:

\begin{lstlisting}[language=Java, caption={Custom Fastify Empty-Body JSON Parser}]
// apps/api/src/server.ts
fastify.addContentTypeParser(
  'application/json',
  { parseAs: 'string' },
  (req, body, done) => {
    if (!body || body.trim() === '') {
      done(null, {});
      return;
    }
    try {
      const json = JSON.parse(body);
      done(null, json);
    } catch (err) {
      done(err as Error, undefined);
    }
  }
);
\end{lstlisting}

\section{Frontend Web Client Implementation (React 18 + Vite)}
The web client, situated in \texttt{apps/web}, is an interactive Single-Page Application (SPA) constructed with **React 18**, **Vite**, and **Tailwind CSS**.

\subsection{Kanban Board and Stage Synchronization}
The core user dashboard presents an interactive Kanban board dividing applications into six operational stages (\texttt{saved}, \texttt{applied}, \texttt{screening}, \texttt{interviewing}, \texttt{offer}, \texttt{rejected}). Users can drag applications across columns or click stage promotion buttons. Every interaction dispatches an asynchronous \texttt{PATCH /api/applications/:id/stage} call that updates local UI state optimistically while handling network rollbacks gracefully upon error.

\subsection{AI Copilot Drawer and Resume Fit Scorer}
An integrated slide-over drawer provides continuous career intelligence:
\begin{itemize}
    \item \textbf{ATS Keyword Analyzer}: Computes dense cosine similarity and lexical keyword overlap between candidate resumes and job requirements.
    \item \textbf{Google XYZ Bullet Point Generator}: A specialized rewriting assistant that prompts the user for rough job achievements and reformulates them into impactful resume bullets adhering to:
    \begin{quote}
        \textit{``Accomplished [X] as measured by [Y] by doing [Z]''}
    \end{quote}
\end{itemize}

\subsection{Recruiter Review Panel with Diff Preview}
When inbound emails are synchronized from Gmail, they populate a ``Mailbox Review'' panel. The client displays extracted entity metadata (Company Name, Detected Role, Extracted Intent) side-by-side with current database records, enabling users to accept or dismiss proposed stage changes with a single click.

\section{AI Microservice Implementation (FastAPI + PyTorch)}
The intelligence layer resides in \texttt{services/ai}, exposing an asynchronous REST interface via FastAPI.

\subsection{Zero-Shot Named Entity Recognition (GLiNER)}
To extract unstructured recruiter intent without manual pattern definitions, the service loads \texttt{urchade/gliner\_small-v2.1}. Incoming email bodies are tokenized and processed across dynamic entity labels:
\begin{lstlisting}[language=Python, caption={Zero-Shot Entity Extraction with GLiNER}]
# services/ai/app/extraction.py
labels = ["company name", "job title", "interview date", "rejection", "offer"]
entities = gliner_model.predict_entities(
    email_text, 
    labels, 
    threshold=0.45
)
\end{lstlisting}
Extracted spans are cross-referenced with user application records to identify candidate matches with high confidence.

\subsection{Dense Embedding Provider (BGE-Small)}
The service instantiates \texttt{BAAI/bge-small-en-v1.5} on the CPU. Input titles and description snippets are pre-pended with the standard query instruction prompt and encoded into normalized 384-dimensional vectors. Batch vector dot-products compute cosine similarities across blocked candidates in sub-10ms intervals.

\subsection{Dual-Engine Fallback Resilience}
To guarantee that the platform remains completely functional even in lightweight or offline environments lacking PyTorch neural runtimes, the AI microservice implements an automatic **Dual-Engine Fallback Architecture**:
\begin{itemize}
    \item \textbf{Primary Mode}: Neural execution utilizing PyTorch, Transformers, and GLiNER.
    \item \textbf{Fallback Mode}: If model weights are unavailable or CPU memory is constrained, the service automatically switches to high-speed deterministic heuristics: regex-based company entity parsing, RapidFuzz lexical matching, and token-overlap keyword scoring, guaranteeing 100\% system availability.
\end{itemize}

\section{Asynchronous Background Queue and Outbox Pattern}
Distributed email ingestion and long-running job portal synchronizations are decoupled from HTTP request cycles through a transactional **Outbox Pattern** implemented directly within PostgreSQL.

\subsection{Concurrency Management with SKIP LOCKED}
Background queue workers execute every 120 seconds. To prevent race conditions across parallel worker instances without introducing external dependencies such as Redis or RabbitMQ, task selection utilizes PostgreSQL row-level locking:

\begin{lstlisting}[language=SQL, caption={PostgreSQL Transactional Outbox Worker Query}]
SELECT id, task_type, payload, attempts
FROM background_tasks
WHERE status = 'queued' AND (scheduled_at <= NOW() OR scheduled_at IS NULL)
ORDER BY scheduled_at ASC
FOR UPDATE SKIP LOCKED
LIMIT 10;
\end{lstlisting}

This guarantees that:
\begin{enumerate}
    \item Each task is exclusively processed by exactly one worker instance.
    \item Concurrent workers immediately skip locked rows without blocking.
    \item Tasks failing during execution are retried with exponential backoff up to a maximum threshold before transitioning to \texttt{failed}.
\end{enumerate}

\section{Cloud Deployment and Serverless Infrastructure}
The production deployment of Track is hosted entirely on **Amazon Web Services (AWS)** in the Asia Pacific (Mumbai) region (\texttt{ap-south-1}), orchestrated through Serverless Framework v4.

\begin{figure}[htbp]
\centering
\begin{tikzpicture}[
    node distance=1.5cm,
    cloud/.style={rectangle, draw=black!60, fill=yellow!10, dashed, thick, rounded corners, minimum width=11cm, minimum height=6cm, align=center},
    item/.style={rectangle, draw=blue!70, fill=blue!10, thick, rounded corners, minimum width=2.8cm, minimum height=1cm, align=center, font=\scriptsize},
    cfnode/.style={rectangle, draw=orange!80, fill=orange!10, thick, rounded corners, minimum width=2.8cm, minimum height=1cm, align=center, font=\scriptsize},
    arr/.style={draw, -latex', thick}
]
    % AWS Cloud Boundary
    \node [cloud] (aws) at (4.5, 0) {};
    \node [above right, font=\footnotesize\bfseries] at (aws.north west) {AWS Cloud Infrastructure (ap-south-1 / us-east-1)};

    % Nodes
    \node [cfnode] (cf) at (0.5, 0.5) {Amazon CloudFront\\CDN Distribution\\(track.mourya.tech)};
    \node [item] (s3) at (4.5, 1.8) {Amazon S3 Bucket\\(track-platform-web)};
    \node [item] (apigw) at (4.5, -0.8) {HTTP API Gateway v2\\(Regional Proxy)};
    \node [item] (lambda) at (8.5, -0.8) {AWS Lambda\\(Node.js 20 ARM64)};
    \node [item] (eventbridge) at (8.5, 1.8) {EventBridge Rule\\(cron(*/2 * * * ? *))};

    % Connections
    \draw [arr] (cf) -- node[above, font=\tiny]{Default (*)} (s3);
    \draw [arr] (cf) -- node[below, font=\tiny]{/api/*} (apigw);
    \draw [arr] (apigw) -- node[above, font=\tiny]{Invoke} (lambda);
    \draw [arr] (eventbridge) -- node[right, font=\tiny]{Queue Trigger} (lambda);
\end{tikzpicture}
\caption{Production Serverless Cloud Topology on Amazon Web Services}
\label{fig:serverless_infra}
\end{figure}

Key infrastructure configurations include:
\begin{itemize}
    \item \textbf{Single-Bundle ES Lambda}: Fastify backend code and dependencies are bundled via \texttt{esbuild} into a single 3.5MB file (\texttt{dist/lambda.mjs}) with custom \texttt{\_\_dirname} ESM shims.
    \item \textbf{Cost-Optimized Architecture}: Running on ARM64 Graviton processors maximizes cost-performance ratios while remaining permanently within the AWS Free Tier.
    \item \textbf{Edge Distribution \& SSL}: CloudFront caches static SPA assets globally while routing API calls dynamically to API Gateway. SSL certificates are provisioned via AWS Certificate Manager in \texttt{us-east-1} and verified via Cloudflare DNS.
\end{itemize}

